\section{10. Limitations}

\subsection{10.1 Not a Compression Method}

ClusterFormer \textbf{increases} parameter count by ~{{PARAM\_CHANGE}}\% due to 
LoRA. It is a \textbf{compute efficiency} method, not a size reduction.

\subsection{10.2 Short-Sequence Overhead}

For $L < 128$, ClusterFormer is slower than baseline. The clustering 
overhead (K-means iterations) exceeds the savings.

\subsection{10.3 Decoder Path Uses Exact Attention}

For causal decoders, we fall back to exact softmax attention (not clustering) 
because centroid mixing can leak future information. This limits 
ClusterFormer's benefit on autoregressive generation.

\subsection{10.4 Clustering Quality}

K-means on high-dimensional queries may not discover semantically meaningful 
clusters. The clustering is a \textbf{heuristic}, not a principled attention 
approximation.

\subsection{10.5 No Formal Approximation Bound}

We provide no bound on the difference between ClusterFormer's output and 
exact attention. The approximation error depends on K-means distortion 
and clustering quality.

\subsection{10.6 Task Coverage}

Only SST-2. Generation, QA, NLI, and long-context tasks are untested.

\subsection{10.7 Model Coverage}

Three models (66.9M--124.6M). Larger models untested.

\subsection{10.8 Hardware}

Only RTX 4070 Laptop. Scaling to A100/H100 unknown.

\subsection{10.9 Learned Hyperparameters}

Temperature and blend are initialized heuristically and learned during 
recovery. Their optimal values may depend on the task.

\subsection{10.10 Single Seed}

Standard error ~1.69\%. Differences below 3.4\% may be noise.

\subsection{10.11 FlashAttention Compatibility}

ClusterFormer does not use FlashAttention kernels. A fused kernel could 
provide further speedup but is not implemented.

\subsection{10.12 Code Quality}

Research-grade:
\begin{itemize}
\item No unit tests for edge cases (empty clusters, single-token sequences)
\end{itemize}
\begin{itemize}
\item Broad exception handling
\end{itemize}
\begin{itemize}
\item K-means iterations hardcoded
\end{itemize}

\subsection{10.13 Ethical Considerations}

Dual-use. See Appendix E.